\chapter{演化悖论：技术加速与物理约束的鸿沟}

打开手机，屏幕里的内容一天换一轮；推开窗户，楼下的街道十年不动。暴雨之后的老旧管网照样积水，早高峰的高架桥照样堵，一份外卖照样要一个骑手顶着雨骑上半小时。生成式模型可以在六十秒内算出一段以假乱真的街头视频，却没法替人下楼取一次快递。两种速度都在跑，问题出在它们常常被当成同一个速度。

\section{技术加速的折叠幻觉与黏滞状态}

\subsection{认知的错位}

屏幕里的内容更新得越快，人越容易把更新的速度当成改变的速度。就业、教育、产业和政策的判断因此常把两个过程混在一起：信息在表示空间里被生成，材料、能量和生命在现实中被重新安排。前者可以依靠复制和重试迅速扩张，后者还要穿过执行、验证和协调。两种时钟被强行对齐，阻力便从判断中消失了。

\subhead{两种速度}

生成式人工智能带来的震动，来自写作、绘图、编程和规划的速度变化。一份方案可以在几秒钟内出现，人们便容易把这种变化外推到方案所描述的世界。比特速度与现实速度在这里被混为一谈。前者显示候选出现得多快，后者还取决于候选能否穿过具体系统。

AI 确实极大降低了搜索、生成、模拟和表达成本。它可以快速提出建筑方案、交通模型、药物候选、商业计划、代码和政策文本。但现实不是一段可以直接改写的文本，也不是一张可以随时重绘的图片。现实由材料、能量、身体、时间、设备、制度、责任和复杂反馈构成。一个方案被生成，不等于它已经被执行；一个模型给出答案，不等于现实系统已经完成状态转移。AI 可以让“找到可能解”变得便宜，却不能自动让“把解变成现实”同等便宜。换言之，AI 的本质能力是更快地发现、生成和筛选路径，而不是取消路径被执行、验证和承担的现实代价。

这种错位有一种通俗的测量方式。大众认知的时钟被拨到了 22 世纪，物理世界的时钟还停在牛顿力学的轨道上，某些基础设施领域，比如老旧的电网和排水系统，甚至保持着 19 世纪的节奏。心理预期与物理现实之间的裂痕没有因为算法的进步而弥合，反而在拉大。差距本身不产生危险，危险来自对它失去感知，而失去感知的第一步，是把屏幕里的经验直接当成世界的经验。要看清这个误读是怎么发生的，需要分别看两边。

数字表示具有低复制成本、低传播成本和便于恢复的工作状态。模型训练完成后，新增一次生成主要消耗推理所需的计算资源；草稿和测试版本可以被替换，代码也可以退回先前保存的状态。使用者因此获得了快速重试的条件。这种便利取决于操作尚未触发外部后果：已经泄露的信息、执行的交易和进入生产系统的错误，都有各自的恢复成本。

具体的例子比属性更有说服力。一段光影正确、物理看似合理的城市街景视频，可以在六十秒内被算出来，而它调动的只是电子在硅晶片上的微观跃迁。一次翻译、一次摘要、一次代码补全，背后是同样的跃迁，代价按厘计费。于是在很短时间内，写作、翻译、摘要、搜索、客服、编程和图像生成这些活动，被从“需要训练的能力”改写成“需要调用的接口”。人感受到的速度提升是真实的，问题在于，被提升的只是这条链上的第一段。

两个域的差别可以写成两个指数项不同的增长律，

\begin{equation}
u_{\text{比特}}(t)=u_{\text{比特}}(t_{0})\,e^{a(t-t_{0})},
\qquad
u_{\text{原子}}(t)=u_{\text{原子}}(t_{0})\,e^{b(t-t_{0})}
\end{equation}

式中$u_{\text{比特}}$与$u_{\text{原子}}$分别代表选定任务在单位成本下的操作量，$t_0$为比较起点，$a$、$b$为给定区间内拟合的增长率。若$a>b$，两者归一化比值的对数以$a-b$增长。这是一种描述速度分离的工作形式。算力、能源转换和物质吞吐各有统计口径，增长率须由各自数据确定；图\ref{fig:two-speeds}用示意曲线说明差别。具体系统可能呈现指数、线性、阶跃或饱和增长，曲线形状由技术与过程共同决定。

\begin{figure}[ht]
\centering
% >>>>>>>>>> figs/fig-two-speeds.tex >>>>>>>>>>
% 两种增长速度的概念示意，不代表历史测量数据。
\begin{tikzpicture}[x=1cm,y=1cm,font=\small]
\useasboundingbox (-0.25,0.85) rectangle (12.95,-6.35);
\draw[-{Stealth[length=2mm]},ink] (1.1,-4.5)--(12.0,-4.5) node[above left,font=\scriptsize]{时间（示意）};
\draw[-{Stealth[length=2mm]},ink] (1.1,-4.5)--(1.1,0.35) node[above right,font=\scriptsize]{归一化能力（示意）};
\fill[slate!11] (1.1,-4.35)--(3.0,-3.9)--(5.0,-3.15)--(7.0,-2.1)--(9.0,-0.9)--(11.0,0.1)--(11.0,-3.2)--(9.0,-3.5)--(7.0,-3.75)--(5.0,-4.0)--(3.0,-4.2)--cycle;
\draw[deepblue,line width=1.8pt] (1.1,-4.35)--(3.0,-3.9)--(5.0,-3.15)--(7.0,-2.1)--(9.0,-0.9)--(11.0,0.1);
\draw[terra,line width=1.8pt] (1.1,-4.35)--(3.0,-4.2)--(5.0,-4.0)--(7.0,-3.75)--(9.0,-3.5)--(11.0,-3.2);
\node[anchor=west,text=deepblue,font=\scriptsize] at (4.2,-0.8){增长较快的表示能力};
\node[anchor=west,text=terra,font=\scriptsize] at (6.0,-4.12){受过程与产能约束的执行能力};
\draw[{Stealth[length=1.6mm]}-{Stealth[length=1.6mm]},ink!65] (11.3,-3.2)--(11.3,0.1);
\node[anchor=west,font=\scriptsize,align=left] at (11.4,-1.5){速度\\分离};
\node[anchor=north west,font=\scriptsize,text=slate,align=left] at (1.4,-5.15){固定任务与成本口径后，分别观察两条能力曲线。\\图形说明可能的分离，不指定历史增长率。};
\end{tikzpicture}

% <<<<<<<<<< figs/fig-two-speeds.tex <<<<<<<<<<

\caption{两种速度的概念示意。横轴与纵轴采用示意刻度；曲线说明表示能力与执行能力可能以不同速度改善。具体增长率须按固定任务和成本口径测量，曲线间距不等于广义黏滞系数。}
\label{fig:two-speeds}
\end{figure}


\subhead{窗外的 20 世纪：被质量与熵增锁死}

这种能力进入物质过程时，要重新支付执行的代价。视频中的交通流由计算生成，现实中的车辆则要由能源驱动，经过制动、转向和接触控制。跨海大桥的图纸可以快速修改，施工仍要组织材料、设备和人员。钢材强度、混凝土强度发展和地基条件进入的是同一个工程账本。混凝土的28天强度是常用检验指标，配合比、温度和养护方式可以改变强度发展的速率；图纸生成得更快，只改变了这条链的设计环节。

第二重约束是存量。城市建立在过去规划的道路网之上，地下管线穿过已经有人生活的街区。自动驾驶和城市智能系统需要多少设施改造，取决于技术路线；有的尽量复用道路，有的增加感知、通信和电力设施。无论采用哪一条路线，车辆、道路使用者、维护体系和责任安排都要衔接。不能像删除代码一样删除一条老旧街道，因为那里住着真实的人，流着真实的物资。现实中的更新，首先是一场保持系统继续运转的切换。

第三重约束是熵。模拟器的规则和边界由设计者规定；现实环境还包含未建模的扰动、状态变化和参与者。AI 可以在模拟器里跑出一亿公里无事故的驾驶记录，现实中一个乱穿马路的行人、一场突如其来的暴雨、一块被遮挡的路牌，都可能让系统停下。物理世界的混乱程度远超数字逻辑的覆盖范围。修路队面对的也不只是铺路，还有地质沉降、管线冲突、周边居民投诉和行政协调。算法可以模拟秩序，无法瞬间消弭现实中的熵增。

\subhead{模拟与现实的范畴错误}

三个属性当中，可逆性最容易被误用。数字孪生中的时间有某种可逆性，回滚一次只需要一条命令，这鼓励了激进的迭代，因为失败的惩罚趋近于零。原子世界没有这条退路。拆掉的桥墩不能回滚，误服的药不能撤销，报废的电池不能重生成。可逆与不可逆的差别不是程度差别，而是类型差别：它决定了试错成本函数长什么样。在可逆系统里长出来的一套方法论，直接搬到不可逆系统里，得到的往往不是更快，而是更贵。

更隐蔽的一步，是把预测能力直接当成干预能力。模型可以从数据中学到杯子下落与破碎之间的关系，机器人要接住杯子，却还需要估计瞬时位置、输出合适的力矩、建立接触，并根据反馈修正动作。学习模型可以参与每一步，前一步的成绩仍然不能替后一步交卷。生成视频越流畅，窗外的施工队就越显得笨拙；这种比较漏掉了施工队真正需要完成的大部分工作。

公共讨论里的“接管”推论，经常把这段距离省掉：模拟已经逼真，执行便只剩时间问题。实际需要跨过的是传感、执行器、能量、接触与反馈。像素收敛了，力矩还要重新控制；图纸通过了，构件还要经受载荷。问题的关键，正是这两张成绩单之间的空白。

因此这一节的核心判断可以压成一句话：\key{在表示空间里成立的速度，到了现实空间要重新计账。}物理世界有很高的黏滞，它由能源成本、材料属性、物流半径和社会协作的复杂度共同构成。这层黏滞吸收了 AI 指数级增长的冲击波，也构成了对抗技术狂热的第一道防线，它客观上延缓了算法对实体经济的冲刷速度，为人类留出了喘息和转型的时间。这不是坏事，但它有一个容易滑过去的前提：黏滞只有在被承认的时候才是缓冲，如果它被当作不存在的噪声，人就不会为它留出任何冗余。

到这里可以回答开头那个问题。屏幕里的内容一天换一轮，街道十年不动，两种速度的差距本身就是事实，不需要额外的解释。需要解释的是另一件事：为什么在差距如此明显的情况下，公共判断仍然持续偏向快的那一侧。答案不在物理世界，而在信息如何被筛选、放大和定价。这个机制有一个已经用了很多年的名字。

\subsection{媒体的幸存者偏差}

\subhead{注意力如何被筛选}

这种错位有它的来源。最轻、最快、最适合演示的技术，占据了最多注意力。视频生成可以当场展示，给服务器供电的变电站却很少进入镜头。芯片制造、电力传输、冷却、采矿和施工各有自己的进步速度，不能从聊天窗口的更新频率推出来。被展示的是界面，被省略的是支撑界面的世界。

筛选机制并不神秘。媒体和资本市场奖励的是能在短时间内被展示、被复制、被估值的对象；一座数据中心的混凝土浇筑过程没有传播价值，一段能生成视频的模型有。结果是可展示的部分获得了完整的叙事，不可展示的部分连名字都没有。这种不对称会自我强化：越是被讲述的东西越容易获得资源和关注，越容易获得资源的东西越容易被继续讲述。几十年下来，公众拿到的样本已经被筛选过一遍，而人用这个样本来推断整体的性质。

偏差还有一个心理学版本。人脑对指数增长的曲线有天然的偏好，对线性增长的曲线缺乏耐心。软件故事讲的正好是指数曲线：用户数在几个月里翻十倍，估值在两年里翻几十倍。物理世界的故事是线性曲线：基站一座一座建，光缆一米一米铺，良率一个点一个点爬。前者适合传播，后者只适合执行。于是归纳谬误出现了：我们误以为整个世界都已经轻量化了。

这个谬误一旦成立，推论就走得很远。所有产业，无论医疗、教育、能源还是制造，都被期待遵循摩尔定律的曲线跃升；任何行业的落后，都被归因于不够数字化或者 AI 不够聪明。它会导向一个具体的预判：只要模型的大脑足够强，它就能像接管聊天窗口一样接管每一条流水线和每一辆重型卡车。人们记住了代码的轻盈，遗忘了重力的刚性。而这类误判最难纠正的地方在于，它在被证伪之前，一直是自我印证着的：被筛选掉的样本从不开口。

\subhead{沉默的基座：遵循热力学的原子世界}

被这套叙事折叠掉的部分要庞大得多。能源开采、重型物流、基础设施建设、精密制造、农业生产，这些不在头条里的产业构成全球实体经济的绝大部分，它们遵循的不是摩尔定律，而是热力学定律。在这个角落里，效率提升不靠复制粘贴，靠对物理世界硬碰硬的改造；没有零边际成本的奇迹，只有每一份做功都必须支付能量的铁律。误判之所以能持续，是因为这一侧的反馈很慢，慢到不足以打断另一侧的叙事。

通信行业是最清楚的例子。手机上的信号格给人一种瞬时覆盖的印象，背后却有管道、光缆、机房、铁塔与检修。铺设光纤要协调道路和地下管线，设备升级要安排供电、运输、安装和测试。协议更新与实体网络更新，关键路径各不相同。差别来自要完成的操作，而不是某个“原子天然比电子慢”的普遍倍数。

能源产业把这段距离放得更大。调度算法可以改善既有容量的使用，新增容量仍要经过设备、土地、融资、施工与接网。伯克利实验室2026年版《Queued Up》统计，截至2025年底，美国约有2060吉瓦发电和储能容量在排队接入；在有数据的地区，2025年建成项目从申请接网到商业运行的中位时长超过五年。\cite{berkeley2026queues} 队列里包含最终会退出的项目，排队容量不能直接算成未来供给。一个电源被设计出来、一份接网申请被提交、电力真正送到负荷端，是三件事。

精密制造的对比更加刺眼。光刻设备内部凝结着高度复杂的知识，制造、装配、运输、安装与校准却一步也不能停在知识里。良率爬坡来自工程师与材料误差、工艺波动之间的持续交手。算法可以缩短反馈周期，也可以让每轮试验更有价值；设备仍要经历这一轮反馈。

\subhead{云端的真相：地球上最沉重的实体}

幸存者偏差最极致的体现，是对“云”这个词的集体误读。在公众想象里，云是轻盈的、去物质化的：数据上传之后似乎脱离了地球引力，进入一个纯粹的比特天堂，似乎只要生活在云端，就不再受物理资源的限制。从物理上看，云很可能是地球上最重、最依赖物理资源的实体之一。所谓的云端不在天上，它躺在地上，是坐落在荒漠、深山或者海底的超大型数据中心，由钢筋混凝土包裹，为周边的供电、散热和安防配套一整套工程。

算力的物理底座，如今已经挤进了电网规划。国际能源署2026年10月14日发布的《Powering Artificial Intelligence》把可用网络容量、接网队列和许可周期列为数据中心建设的约束。\cite{iea2026powering} 芯片、机架、电力设备、铜缆与光纤，各有供应链；冷却方式、气候和负载共同决定水与电的需求。云端每增加一份能力，地面就要回答它放在哪里、怎样供电、热排到哪里。

回看传播规律，这套偏差并不需要谁去刻意制造。可展示的部分自动获得注意力，不可展示的部分自动失去注意力，两者叠加就得到了一个重量被系统性低估的世界图景。这也解释了一个反常现象：关于 AI 的讨论越热烈，关于电力、冷却、矿产和供应链的讨论反而越少，因为后者不进镜头。要给这个偏差纠偏，需要的不是更多报道，而是一个能同时容纳两边的概念。

\subhead{给阻力起一个名字：技术黏滞}

本章用\term{技术黏滞}描述这种差异。所谓技术黏滞，是指智能、算法或技术方案推动现实系统从状态 A 转移到状态 B 时，遭遇的广义阻尼。它不是简单地说现实变化缓慢，而是强调：现实改变不只取决于是否存在一个解，还取决于这个解能否被执行、被验证、被维护、被协调，并在失败时被承担后果。AI 的核心优势在于降低搜索成本，但现实状态转移还要支付执行成本、验证成本、维护成本、协调成本和责任成本。这里的关键词是“广义”。黏滞在物理学里指流体内部阻碍相对运动的性质，本书借用的正是这层含义：阻力不来自某个外部障碍，而来自介质本身的内部结构。现实系统的介质属性，才是智能进入时真正遇到的东西。

由此可以推出一个反直觉但重要的推论：\key{生成越便宜，可信结果反而越稀缺；方案越丰富，现实吸收能力越成为瓶颈。}在稀缺结构发生变化的地方，竞争结构也会变化。过去，许多领域稀缺的是生成能力：能不能写出文本，能不能画出图像，能不能找到分子，能不能提出方案，能不能生成代码。现在，生成能力迅速膨胀，候选方案急剧增加，信息空间开始出现通胀。问题不再是没有方案，而是方案太多；不再是没有答案，而是答案太便宜；不再是缺少表达，而是缺少可信表达。未来的瓶颈将越来越从生成转向验证，从创造转向确证，从可能性转向可承担性。

任何现实状态转移，都至少要穿过三类约束。要有足够资源推动，要有足够时间展开，要能在既有结构中传导。为了分析这种阻尼，可以提出一个最小模型：RTG，即\term{资源势垒}（$R$）、\term{时间刚性}（$T$）和\term{拓扑耦合}（$G$）。资源势垒回答一个改变是否推得动，时间刚性回答一个改变是否来得及，拓扑耦合回答局部改变能否传导为系统改变。三者共同决定智能进入现实系统时的黏滞程度，这三类约束也构成理解全书的最小框架。

\subsection{虚拟逃逸论的物质依附}

\subhead{向内坍缩}

既然现实这一侧又重又慢，最省力的反应就是绕开它。于是产生了一种逃逸倾向：面对沉重的物理世界，技术想象不断向去实体化方向移动，试图在纯符号层面寻找自由。这种逃逸有一个内在悖论，虚拟世界越繁荣，对高能耗物理基座的依赖越深。试图脱离引力，却必须消耗更多燃料。

这种倾向在最近十年获得了一个更完整的形式，可以称为虚拟逃逸论。它的前提是一组真实的挫败感：摩天大楼的房价高不可攀，城市交通的拥堵积重难返，通过物理迁移改变处境的通道在收窄。一个文明能够支配的功率总量是有限的，物理世界的重构因此既慢又贵。既然在原子世界里争取一个物理坐标的成本如此之高，不如把生活的重心整体移到比特世界。它在更深一层呼应着一种古老的冲动：通过纯粹的知识摆脱物质枷锁，把肉身当作需要卸下的负担。硅谷在这里扮演了新的布道者角色，代码被许诺为通向自由形态的路径。

核心变化是价值排序。实体世界在逃逸者眼中正在失去吸引力，甚至变成需要被超越的累赘。传统的衣食住行被降格为维持生命体征的底线需求，一种肉体极简、灵魂极繁的生存策略开始流行。戴上头显设备就能置身海滩，为什么还要背三十年的房贷去争一个市中心的位置；工作可以远程完成，社交可以在虚拟空间进行，聚会可以在数字场景里举办，那么大排量汽车和早高峰拥堵就失去了意义。这是一种主动的空间退却：把生存空间从广袤的地理坐标，压缩到一个光纤接口上。

与这种生活方式配套的，是一种政治想象。以 Balaji Srinivasan 为代表的一批硅谷思想家提出过“网络国家”的概念，主张云端先行、落地在后，用软件的灵活性规避传统国家的僵化。从物理第一性原理看，这类构想有一个共同的漏洞：它们寄生在既有国家的物理基座之上。分散在巴厘岛或者里斯本的数字游民既不生产电力，也不铺设海底光缆，更不具备维护航道安全或者处理大规模传染病的能力。他们依附于那个由钢铁和化石能源构成的、被他们试图逃离的母体文明。任何无法解决独立能源与物质闭环的数字乌托邦，只能是依附于实体经济的俱乐部，而不是一种独立的文明形态。

\subhead{热力学的账单}

这套构想包含一个被技术光环掩盖的物理悖论。虚拟世界并不独立存在，它只是物理世界的一种排列形式：在微观层面，它是硅晶片上受控运动的电子，是硬盘磁头上精准的翻转，是光纤中全反射的光信号。维持一个高度有序的虚拟世界，必须以物理世界的高熵增和大量能量耗散为代价。这是热力学第二定律开出的账单，没有逃单的路径。

每一次推理、渲染和通信，都由具体的电路与设备完成。电力来自哪里的哪一种电源，取决于地点、时段和采购安排。把活动移到线上，可能节省交通和空间，也会增加计算、通信与设备需求。净效应要把被替代的活动和新建的设施一起核算。虚拟世界可以改变账单，无法脱离开账单的地方。

去空间化的生存，依赖极高成本的空间基础设施。一个人可以躲进深山当数字游民，前提是深山里必须有稳定的电力供应、高带宽的通信基站，以及能在一小时内送达急救药物的物流网络。云端不是飘在天上的云，而是躺在地上、由数万吨钢筋混凝土包裹的超大型数据中心；这些建筑需要防震、防火、防潮，需要庞大的冷却系统带走算力排出的废热。服务器、芯片、光纤、铜缆全部是物理实体，它们的开采、冶炼、制造和运输遵循最传统的重工业逻辑。虚拟世界的每一秒轻盈，都以原子世界更深重的负重为代价。

数字系统的连续运行依赖一组可失效的条件。海底光缆、供电、底层依赖与维护团队，各自连接不同的故障路径。系统可以通过冗余、异地恢复和隔离提高可靠性，也会为此增加资源支出。维护成本与鲁棒性需要按具体架构比较，热力学并不规定有序度越高就越脆弱。虚拟生活的物理条件，需要在这些设施与工作中核算。

还有一笔账被转移到了地理上。界面把结果留在眼前，把供电、散热、维修和废弃物处理放到别处。屏幕上的整洁不是热力学熵；要说明物理代价，需要考察执行这些操作的设备与环境。对一个与环境合起来构成孤立整体的系统，第二定律写成，

\begin{equation}
\Delta S_{\text{局部}}+\Delta S_{\text{环境}}=\Delta S_{\text{总}}\ge 0
\end{equation}

式中三个熵变的单位均为焦耳每开尔文。若局部熵减少，环境熵至少增加同样大小，等号对应理想可逆过程。电力转换、冷却、制造和运输包含不可逆过程，真实成本还取决于具体路线。这个总账约束物理过程，不能给每张图像、每份知识或每一寸整洁屏幕标出固定的熵价。水耗、污染与电子废物也要分别测量。

\subhead{生命的锚点：不可上传的肉身}

除能源与基建之外，还有一个更基础的约束：生命本身。只要意识还锚定在神经网络中无法上传，只要生老病死这套生物过程还没有被破解，人就没有办法彻底离开这个黏滞的世界。肉身受制于极严苛的物理常数：它需要 37 摄氏度的恒温，需要精准的生化反馈，会因为几摄氏度的体温波动或者微量细菌入侵而崩溃。这种生物性的脆弱是演化写定的。软件可以瞬间升级，伤口不能瞬间愈合，大脑发育的周期也不能压缩。这种时间上的不可压缩性，构成生命最本质的黏滞。

所谓轻量化的生活，实质上是一场关于秩序与混乱的置换。当一个人沉浸在零摩擦的虚拟体验中，这种顺滑并非凭空产生，它通过复杂的社会分工，把物理负担转嫁给了底层的原子系统。数字生活的每一次轻盈，都在物理底层结算为更深的磨损：维护服务器的电工、风雨中的骑手、排出废热的冷却塔。虚拟层面的每一寸整洁，最终都由物理世界的耗散来代偿。

因此不存在所谓的逃逸。虚拟逃逸论是面对物理壁垒时的一种认知撤退，它试图在比特世界寻找廉价捷径，结果陷入更深的基质依赖。这揭示了一个反向关系：\key{屏幕里的每一寸整洁，都在屏幕外结算为耗散。}越是想在比特世界飞升，就越需要往原子世界扎根。没有热力学底座，就没有虚拟的幻梦；没有健康的生物载体，就没有意识的栖息。那些看上去过时的重资产行业，修路、采矿、发电、种粮，因为它们支撑着数字世界的运转，反而成为文明中最难撼动的一部分。

这也把问题推向一个更基本的位置。物理约束在这里一直以阻碍的形式出现：它让方案落不了地，让逃逸者被拽回来，让加速的承诺打折。但如果把观察的尺度放大，同一批约束还有另一副面孔。它们不只是让变化变慢的东西，也是让系统不至于散架的东西。要看清这一面，需要把黏滞从阻力改读为阻尼。
